\section{Introduction}\label{sec:intro}

% DRAFT: to be revised after the full texts of [Wang2026PretrainFEM], [Sun2026FEVessel]
% and [MonteiroFernandes2025] have been read (outline v3, Sec. 7).

Neural surrogates of finite-element analyses are trained to return, in one network
evaluation, the field that an assembly and solve would return, so that design studies,
optimisation and uncertainty quantification can query many configurations cheaply
\cite{Pfaff2021,Lu2021DeepONet,Li2021FNO,He2024GeomDeepONet,Zeng2025PCNO,Liu2026GINOT}.
Most are trained by supervised regression on reference solutions. Each reference
solution is a full solve of the system that the surrogate is meant to replace, so the
training set costs, before the surrogate has made its first prediction, as many solves as
it has instances and load cases; every extension of scope, to a new geometry family, load
regime or mesh resolution, costs more.

For linear elastostatics there is a training signal that needs no solve. The
finite-element solution minimises the assembled discrete potential energy
$\Pi_h(u)=\tfrac12u^{\top}Ku-F^{\top}u$, whose stiffness matrix $K$ and load vector $F$
every finite-element code assembles before it solves; evaluating $\Pi_h$ at a prediction
$u$ costs one sparse matrix-vector product. The energy of a prediction exceeds that of
the solution $\Ustar$ by half the squared stiffness-norm error,
$\Pi_h(u)-\Pi_h(\Ustar)=\tfrac12\norm{u-\Ustar}_K^2$, a classical identity of the
energy-norm theory of finite elements \cite{Ciarlet1978}, and this half squared norm is the
strain energy of the error's stress field. Minimising the energy over a set of training
instances is therefore exactly supervised regression in the stiffness norm, with the same
gradients at every point: a surrogate trained on the energy learns from a weighted
mean-square stress error, and no training instance is ever solved.

Energy minimisation as a training principle is not new. The deep Ritz method
\cite{EYu2018} and the deep energy method \cite{Samaniego2020} minimise a continuous energy
over neural trial functions, and later work trains operators and mesh-based networks
without labels, with residual or energy losses
\cite{Xu2024VOL,Eshaghi2025VINO,Lee2025FEONet,Dalton2023,Wurth2024,Wang2026PretrainFEM}
(\cref{sec:related}). That the energy ranks candidate fields is the Ritz minimum principle
itself. This paper adds two things for neural surrogates. The first is the systematic use
of that principle to evaluate and screen their predictions without reference solutions: an
exact link from the energy to the stress error, a test that tells whether a prediction is
worse than the zero field, and a closed-form rescaling that guarantees it is not. The second
is a controlled test of what label-free training buys: the same network trained on the same
instances with and without their solutions, under criteria fixed before the runs, with the
criteria that failed reported alongside those that were met. By the identity above, the
test compares two losses, the stiffness-norm error, which needs no solve, and the usual
relative displacement error, which does; the labels add no information that $K$ and $F$ do
not already carry. The paper extends a companion note \cite{CaoSong2026}, which stated the
exactness identity, the conditioning bounds and their executable checks.

The contributions are as follows.
\begin{enumerate}
\item Theory (\cref{sec:theory}). The label-free energy objective is supervised
stiffness-norm regression (\cref{lem:exact,cor:training}); for isotropic elasticity the
energy gap is a compliance-weighted mean-square error of the stress field and bounds the
error of the von Mises field (\cref{prop:stress}). Displacement and energy-norm errors can
rank two predictions in opposite orders only when their displacement errors differ by less
than the square root of the condition number of $K$, and then they can
(\cref{rem:rank}). The sign of $\Pi_h(u)$ tells, for each load case, whether a prediction is
worse than the zero field (\cref{cor:zero}), and the two terms of $\Pi_h(u)$ give a
rescaling after which it is not (\cref{prop:amp}). Under the classical Chebyshev bound, a
warm start lowers the iteration bound of conjugate gradients by a fraction that depends
only on the prediction's relative energy gap and the target tolerance (\cref{cor:warm}).
\item Experiments (\cref{sec:results2d,sec:results3d}). On three-dimensional tetrahedral
meshes, a transformer trained on the energy alone was more accurate on all five metrics,
in every seed, than the same transformer trained on the solutions of the same 1,024
instances. A graph network trained on those solutions had lower displacement errors than
the label-free transformer on most instances but higher stress errors on almost all of
them, so that displacement error ranked the two in the opposite order to stress error.
Every supervised network returned some in-band predictions worse than the zero field; the
label-free network returned none, averaged over each instance's load cases, and
\numThreeFreeLoadWorse{} in \numThreeFreeLoads{} load cases taken singly. In two dimensions, the label-free network's
displacement error was \numTwoFreeDispExcessAbs{} larger than the supervised network's, and
it was more accurate on every other metric.
\item Negative results (\cref{sec:results2d,sec:transfer,sec:cost}). In two dimensions,
conjugate gradients started from a label-free prediction did not need fewer iterations to a
tight residual tolerance than from zero. No network transferred to a three-dimensional mesh
about twice as fine without retraining; a closed-form rescaling applied after the fact
removed about two-thirds of the label-free network's displacement error there.
\item Cost (\cref{sec:cost}). The label-free network was trained without a single solve;
in three dimensions the supervised networks' 1,024 labelled training instances took
\numThreeTrainSolves{} solves, and evaluating any method on the in-band and fine-mesh
validation sets took a further \numThreeEvalSolves{}.
\end{enumerate}

We do not claim that a surrogate is faster than solving a linear elastostatic problem:
sparse direct and iterative solvers are fast at the sizes studied here
(\cref{sec:cost}). Linear elastostatics is the test bed on which the exactness of the
objective can be proved and the comparison with supervision kept narrow: the same network,
the same instances and the same schedule, with only the loss and its learning rate
differing. The interest of label-free training lies where solves are expensive and labels
scarce, for example in nonlinear and time-dependent problems whose implicit steps are again
energy minimisations (\cref{sec:discussion}).

\Cref{sec:related} reviews related work, \cref{sec:theory} the theory and
\cref{sec:methods} the problems, networks and protocol. \Cref{sec:results2d,sec:results3d}
report the two- and three-dimensional comparisons, \cref{sec:transfer} the transfer to a
finer mesh and \cref{sec:cost} the cost. \Cref{sec:discussion} discusses the results and
their limits. \ref{app:prereg} lists every pre-registered criterion with its outcome and
the deviations from the plans, \ref{app:e1} a negative result on latent regularisation
and \ref{app:repro} how every number is reproduced from the records.
